\documentclass[11pt,reqno]{amsart}

% --- Encoding / fonts ---
\usepackage[T1]{fontenc}
\usepackage[utf8]{inputenc}

% --- Core AMS packages ---
\usepackage{amsmath,amssymb,amsfonts,amsthm,mathtools}

% --- Graphics and lists ---
\usepackage{graphicx}
\usepackage{enumitem}

% --- Page layout ---
\usepackage[a4paper,margin=3cm]{geometry}

% --- Theorem environments ---
\theoremstyle{plain}
\newtheorem{theorem}{Theorem}[section]
\newtheorem{lemma}[theorem]{Lemma}
\newtheorem{proposition}[theorem]{Proposition}
\newtheorem{corollary}[theorem]{Corollary}

\theoremstyle{definition}
\newtheorem{definition}[theorem]{Definition}
\newtheorem{example}[theorem]{Example}
\newtheorem{question}[theorem]{Question}
\newtheorem{conjecture}[theorem]{Conjecture}

\theoremstyle{remark}
\newtheorem{remark}[theorem]{Remark}
\newtheorem{claim}[theorem]{Claim}
\newtheorem{assertion}[theorem]{Assertion}
\newtheorem{notations}[theorem]{Notations}

\newtheorem{assumption}{Assumption}
\newtheorem{hypothesis}{Hypothesis}

\newtheorem*{theoremA}{Theorem A}
\newtheorem*{theoremB}{Theorem B}
\newtheorem*{theoremC}{Theorem C}

% --- Operators ---
\DeclareMathOperator{\IFS}{IFS}
\DeclareMathOperator{\diam}{diam}
\DeclareMathOperator{\vol}{vol}

% --- Hyperref: load late ---
\usepackage[pdfencoding=auto,psdextra]{hyperref}

% --- TOC depth ---
\setcounter{tocdepth}{3}

% --- Title & Authors ---
\title[T]{TransTitle}

\author[A1]{B1}
\address{Department of Pure Mathematics, Ferdowsi University of Mashhad, Mashhad, Iran}
\email{htina@um.ac.ir}

\author[A2]{B2}
\address{Department of Pure Mathematics, Ferdowsi University of Mashhad, Mashhad, Iran}
\email{zaj@ferdowsi.um.ac.ir}

\thanks{Corresponding author: F. H. Ghane (\texttt{htina@um.ac.ir})}

\keywords{Monotone}
\subjclass[2020]{Primary 37D25; Secondary 37A35, 37C40, 37A60.}

\begin{document}

\maketitle


\begin{abstract}
abbbaaaaaaaa
\end{abstract}

\section{Introduction}
\label{sec:intro}


\section{Model}
\label{sec:model}

% Required packages:
% \usepackage{amsmath,amssymb}
% \usepackage{hyperref}

\section{Introduction and model formulation}
\label{sec:introduction}

Understanding the interaction between electrical excitation and
recovery is a central problem in mathematical cardiac
electrophysiology. Reduced models describe this interaction using
a small number of variables, making it possible to investigate
the mechanisms governing resting states, excitation thresholds,
repolarization, and sustained oscillations. Their relatively
simple structure also provides a setting in which qualitative
properties can be established analytically.

Rogers and McCulloch~\cite{RogersMcCulloch1994} introduced a
modification of FitzHugh--Nagumo-type kinetics in a computational
study of cardiac action-potential propagation. Their formulation
couples a voltage-dependent cubic excitation term to a recovery
variable through a bilinear interaction. Although originally
studied in a spatial propagation framework, its reaction terms
also define a two-dimensional ordinary differential equation
describing spatially homogeneous kinetics.

The present work concerns this autonomous planar system.
The purpose is to investigate its qualitative dynamics as an
established cardiac excitation model, rather than to introduce
additional phenomenological interactions.

\subsection{The autonomous Rogers--McCulloch system}
\label{subsec:model}

Let $V_m$ denote the transmembrane potential, and introduce the
normalized voltage
\begin{equation}
\label{eq:normalized-voltage}
u=\frac{V_m-V_r}{V_p-V_r},
\end{equation}
where $V_r$ is a reference resting potential and $V_p>V_r$ is
a reference excited potential. Let $v$ be a dimensionless
recovery variable representing the aggregate processes
responsible for repolarization and recovery of excitability.

After absorbing the constant voltage and capacitance factors
into the coefficients of the voltage equation, the
spatially homogeneous Rogers--McCulloch kinetics can be written as
\begin{equation}
\label{eq:RM-general}
\begin{aligned}
\frac{du}{dt}
&=c_1u(u-a)(1-u)-c_2uv+I(t),\\
\frac{dv}{dt}
&=b(u-dv).
\end{aligned}
\end{equation}
The parameter assumptions are
\begin{equation}
\label{eq:parameters}
0<a<1,
\qquad c_1,c_2,b,d>0.
\end{equation}
Here $t$ denotes time in the chosen physical unit. No
state-dependent change of time is required to obtain
the polynomial vector field.

To study autonomous phase portraits and limit cycles,
we take a constant applied input,
\[
I(t)\equiv I_0,\qquad I_0\geq 0.
\]
Thus, the system under consideration is
\begin{equation}
\label{eq:RM-autonomous}
\boxed{
\begin{aligned}
\dot u
&=c_1u(u-a)(1-u)-c_2uv+I_0,\\
\dot v
&=b(u-dv).
\end{aligned}}
\end{equation}
A dot denotes differentiation with respect to $t$.
System~\eqref{eq:RM-autonomous} is an autonomous planar
polynomial system of degree three. The case $I_0=0$
describes unforced kinetics, whereas $I_0>0$ represents
a sustained depolarizing input.

The coefficients in~\eqref{eq:RM-autonomous} refer to this
normalized formulation. Numerical parameter values taken
from a dimensional formulation must therefore be converted
using the same voltage, current, capacitance, and time
conventions.

\subsection{Interpretation of the variables and parameters}
\label{subsec:interpretation}

The parameter $a$ specifies the intermediate zero of the
cubic excitation term. In voltage coordinates it corresponds
to the reference level
\[
V_{\mathrm{th}}=V_r+a(V_p-V_r).
\]
It is a parameter controlling excitability, rather than a
universal firing threshold independent of the recovery
state and applied input.

The coefficient $c_1$ determines the strength of the intrinsic
voltage kinetics. In the absence of recovery feedback and
applied input, the cubic term is negative for $0<u<a$,
positive for $a<u<1$, and negative for $u>1$. This sign
structure describes a competition between relaxation
toward rest and regenerative excitation.

The coefficient $c_2$ determines the strength of the
voltage--recovery interaction. For $u,v>0$, the term
$-c_2uv$ opposes depolarization. Its dependence on both
variables distinguishes this formulation from the usual
additive recovery term in the standard FitzHugh--Nagumo
equations.

The recovery equation can be expressed as
\begin{equation}
\label{eq:recovery-relaxation}
\dot v=bd\left(\frac{u}{d}-v\right).
\end{equation}
Consequently, for a fixed voltage, the recovery variable
relaxes toward $u/d$ with relaxation time $(bd)^{-1}$.
The parameter $d$ determines the recovery nullcline,
whereas $b$ controls the recovery rate for a fixed value
of $d$. A slow--fast interpretation requires an explicit
comparison of this rate with the voltage kinetics; it
does not follow from positivity of the parameters alone.

Finally, $I_0$ is the applied current expressed in normalized
voltage-rate units. With dimensionless $u$ and $v$, the
parameters $c_1$, $c_2$, and $b$, together with $I_0$,
have units of inverse time, while $a$ and $d$ are
dimensionless.

The variable $v$ is a phenomenological recovery variable.
It should not be identified with a specific ionic
concentration or individual channel-gating probability
without additional justification. Similarly, the model
describes electrical excitation and recovery, rather
than the progression of a particular cardiac disease.

\subsection{Phase space and elementary geometric properties}
\label{subsec:phase-space}

For nonnegative initial conditions, a natural domain is
\[
\mathcal Q=\{(u,v)\in\mathbb R^2:u\geq0,\ v\geq0\}.
\]
Since
\[
\left.\dot u\right|_{u=0}=I_0\geq0,
\qquad
\left.\dot v\right|_{v=0}=bu\geq0,
\]
the quadrant $\mathcal Q$ is positively invariant.
When $I_0=0$, the boundary $u=0$ is invariant;
when $I_0>0$, the vector field points into $\mathcal Q$
along that boundary.

The normalization of voltage does not imply that
$0\leq u\leq1$ is invariant under every applied input.
In particular, at $u=1$,
\[
\dot u=I_0-c_2v,
\]
which can be positive. Physiological interpretation
must therefore account for the voltage range associated
with a chosen parameter set and stimulation protocol.

Forward trajectories in $\mathcal Q$ are bounded.
Indeed, because
\[
c_1u(u-a)(1-u)+I_0\longrightarrow-\infty
\quad\text{as }u\longrightarrow+\infty,
\]
one can choose $U>1$ such that this expression is negative
for every $u\geq U$. Since $-c_2uv\leq0$ in $\mathcal Q$,
\[
u(t)\leq M:=\max\{u(0),U\}.
\]
The recovery equation then gives
\[
v(t)\leq \max\left\{v(0),\frac{M}{d}\right\}.
\]
These bounds also ensure forward completeness for
nonnegative initial conditions.

The equilibria satisfy
\begin{equation}
\label{eq:equilibrium-equation}
v_*=\frac{u_*}{d},
\qquad
c_1u_*(u_*-a)(1-u_*)
-\frac{c_2}{d}u_*^2+I_0=0.
\end{equation}
Thus, there are at most three distinct finite equilibria.
Their number, admissibility, and stability depend on the
parameters.

A global analysis on $\mathcal Q$ concerns the dynamics
of nonnegative states. A classification on the entire
Poincare compactification is a broader mathematical
problem and also includes regions without an immediate
physiological interpretation. These two meanings of
``global phase portrait'' should be distinguished.

\subsection{Related work and research questions}
\label{subsec:related-work}

The original study of Rogers and McCulloch
\cite{RogersMcCulloch1994} emphasized numerical cardiac
impulse propagation and the influence of tissue geometry
and structure. More recently, L\'opez-Zazueta,
Soto-Bajo, and Fraguela Collar
\cite{LopezZazueta2024} investigated numerically the
response of an isolated cell under periodic pulse
stimulation using the Rogers--McCulloch model.

At the analytical level, Hieber et al.~\cite{Hieber2020}
established the existence and uniqueness of strong
time-periodic solutions near stable equilibria for
periodically forced bidomain equations, including
Rogers--McCulloch ionic transport.

These studies motivate, but do not by themselves resolve,
the global classification problem for
\eqref{eq:RM-autonomous}. A periodic response to a
time-dependent stimulus is a different mathematical
object from an isolated periodic orbit of an autonomous
planar vector field. Likewise, a local existence theorem
near a stable equilibrium does not determine the total
number of autonomous limit cycles or the global
configuration of separatrices.

The main questions for the autonomous system are therefore
the following:
\begin{enumerate}
\item
How does parameter space divide according to the number
and stability of equilibria, including nonhyperbolic
configurations?

\item
For which parameter regimes do limit cycles exist or
fail to exist? When they exist, what can be proved about
their number, multiplicity, and stability?

\item
Which local and global bifurcations organize transitions
between the admissible phase portraits, and which
candidate bifurcations are excluded by the parameter
restrictions?

\item
How do separatrices and basin boundaries organize
resting, excited, and oscillatory behavior in
$\mathcal Q$?

\item
If the analysis is extended to the whole plane, which
singularities at infinity occur, and how do their
separatrices connect with the finite dynamics?
\end{enumerate}

These questions specify a research programme rather than
a claim that each problem is absent from the existing
literature. Establishing novelty requires comparison
with previous analyses of the Rogers--McCulloch equations
and with results for broader polynomial or generalized
Li\'enard families that may contain equivalent systems.

In particular, the bilinear recovery coupling prevents
a direct identification with the usual additive-recovery
FitzHugh--Nagumo formulation. However, a different
appearance of the equations is not sufficient evidence
of mathematical novelty: coordinate transformations,
time rescalings, and existing general theorems must also
be examined.

A substantial contribution would consist of rigorous
results beyond numerical examples or a routine local
stability calculation, such as a sharp limit-cycle
bound, a complete bifurcation classification on a
specified parameter domain, or a global description
of separatrix configurations and attraction basins.
Any claim of completeness must specify the phase space,
parameter domain, and treatment of degenerate cases.

Finally, autonomous oscillations in this reduced model
describe repetitive electrical activity at the level
of homogeneous kinetics. Their existence alone does
not establish a mechanism for tissue-level reentry or
a diagnosis of cardiac arrhythmia, which require
additional spatial and physiological information.

\begin{thebibliography}{99}

\bibitem{RogersMcCulloch1994}
J.~M. Rogers and A.~D. McCulloch,
A collocation--Galerkin finite element model of cardiac
action potential propagation,
\emph{IEEE Transactions on Biomedical Engineering}
\textbf{41}(8) (1994), 743--757.
\href{https://doi.org/10.1109/10.310090}
{doi:10.1109/10.310090}.

\bibitem{LopezZazueta2024}
A. L\'opez-Zazueta, M. Soto-Bajo, and A. Fraguela Collar,
A numerical study of a single isolated human ventricular
cell response to a periodic pulse stimulus current using
the Rogers--McCulloch model,
\emph{Communications in Nonlinear Science and
Numerical Simulation}
\textbf{132} (2024), 107912.
\href{https://doi.org/10.1016/j.cnsns.2024.107912}
{doi:10.1016/j.cnsns.2024.107912}.

\bibitem{Hieber2020}
M. Hieber, N. Kajiwara, K. Kress, and P. Tolksdorf,
The periodic version of the Da Prato--Grisvard theorem
and applications to the bidomain equations with
FitzHugh--Nagumo transport,
\emph{Annali di Matematica Pura ed Applicata}
\textbf{199} (2020), 2435--2457.
\href{https://doi.org/10.1007/s10231-020-00975-6}
{doi:10.1007/s10231-020-00975-6}.

\end{thebibliography}










%\section*{Acknowledgments}
%[Funding, collaborations]

\bibliographystyle{plain}
\bibliography{bibliography}
%\printbibliography




\end{document}
